\subsection{Interaction as dependent substitution}
\label{sec:interaction}

\subsubsection{Complete executions}

Fix a closed script, or equivalently instantiate its public context. A profile
is a tuple $\vecp=(p_i)_{i\in\Roles}$ with $p_i:G@i$.
We first define the type of complete executions.

\begin{definition}[Trace and return family]
The trace type is
\begin{equation}
 \Trace(\End\langle R_i\rangle_i)=\Unit,
 \qquad
 \Trace(\Q_r(x:A).H)=\sum_{a:A}\Trace(H[a/x]).
 \label{dep:eq:9}
\end{equation}
For $t:\Trace(G)$, define the terminal return family by
\begin{equation}
 \Ret_i^{\End\langle R_j\rangle_j}(\star)=R_i,
 \qquad
 \Ret_i^{\Q_r(x:A).H}(a,t)=\Ret_i^{H[a/x]}(t).
 \label{dep:eq:10}
\end{equation}
The complete execution type is
\begin{equation}
 \Exec(G):=\sum_{t:\Trace(G)}\prod_{i\in\Roles}\Ret_i^G(t).
 \label{dep:eq:11}
\end{equation}
\end{definition}

The finite product in \eqref{dep:eq:11} is a record indexed by roles. All returned values
refer to the same trace $t$. A separate trace for each endpoint would lose the
agreement that communication is meant to establish.

\begin{definition}[Execution]
\label{def:execution}
Define $\Run_G(\vecp):\Exec(G)$ recursively. At a terminal,
\begin{equation}
 \Run_{\End\langle R_i\rangle_i}(\vecp)
       :=(\star,(p_i)_i).
 \label{dep:eq:12}
\end{equation}
At $G=\Q_r(x:A).H$, destruct $p_r$ as $(a,u_r)$ and put
\begin{equation}
 q_i:=
 \begin{cases}
   u_r,&i=r,\\
   p_i(a),&i\ne r.
 \end{cases}
 \label{dep:eq:13}
\end{equation}
If $\Run_{H[a/x]}(\vec q)=(t,\vec z)$, return
$((a,t),\vec z)$.
\end{definition}

The construction uses dependent elimination of a pair at the owner's endpoint
and dependent application at every other endpoint. It chooses no value on
behalf of the participants.

\begin{lemma}[Dependent interaction]
\label{thm:interaction}
For every well-typed profile $p_i:G@i$, execution is well typed at $\Exec(G)$.
In particular, if $\Run_G(\vecp)=(t,\vec z)$, then
$z_i:\Ret_i^G(t)$ for every $i$.
\end{lemma}
\begin{proof}
At a terminal this is exactly the type of the result record. At a choice,
dependent elimination of $p_r$ gives $a:A$ and $u_r:\Local{r}{H[a/x]}$.
For $i\ne r$, application gives
$p_i(a):\Local{i}{H}[a/x]$, which is definitionally
$\Local{i}{H[a/x]}$ by \cref{lem:substitution}.
The induction hypothesis therefore applies to the common continuation
$H[a/x]$. Its result has the return types required by \eqref{dep:eq:10}, and pairing the
trace with $a$ gives \eqref{dep:eq:11}.
\end{proof}

\subsubsection{The interaction reduction}

A configuration consists of one residual script and a profile of its
endpoints. With canonical head forms, its communication rule is
\begin{equation}
 \Cfg\bigl(\Q_r(x:A).H;
       \Send_r(a,p_r),(\Recv_i(x.q_i))_{i\ne r}\bigr)
 \longrightarrow
 \Cfg\bigl(H[a/x];p_r,(q_i[a/x])_{i\ne r}\bigr).
 \label{dep:eq:14}
\end{equation}
For the reduction semantics, first expand the endpoint macros into $\D$.
Rule \eqref{dep:eq:14} then requires a pair at the owner's head and a lambda at each
receiving head. Label this step by $(r,a)$. Include exactly the following
additional rules: a pure $\D$-reduction of one endpoint term is an internal
configuration step, and a configuration at $\End$ whose endpoint terms are
in data normal form takes one terminal step returning their record.
Terminal configurations have no successors. There are no other steps.
The ordered choice labels reconstruct the trace in \eqref{dep:eq:9}. A complete execution
is a maximal sequence of these steps; pure endpoint reductions may be
interleaved with enabled communication.

\begin{corollary}[Subject reduction for complete connections]
\label{cor:preservation}
A communication step \eqref{dep:eq:14} preserves the typing of all residual endpoints at
one common residual script. Every recorded trace prefix is well typed.
\end{corollary}
\begin{proof}
Invert \eqref{dep:eq:6} and \eqref{dep:eq:7}, and apply \cref{lem:substitution} to every receiver. The
owner's continuation already has the substituted type. The trace clause is
the dependent pair introduction in \eqref{dep:eq:9}.
\end{proof}

\begin{lemma}[Termination and unique execution]
\label{thm:termination}
Assume the normalization, confluence, and canonical-form properties of $\D$.
Every closed well-typed complete connection terminates. It takes exactly
$n(G)$ choice steps and one terminal return. Its complete result is unique up
to data conversion. Interleaving pure endpoint reductions with enabled head
communication does not change that result.
\end{lemma}
\begin{proof}
At a choice, normalization gives the owner's canonical dependent pair. The
receivers are applicable functions at $A$. Subject reduction supplies a
well-typed continuation with one fewer script binder; data substitution leaves
that constructor count unchanged. Induction on $n(G)$ gives termination and
the step count. Each pure normalization phase terminates by the assumption on
$\D$. There are finitely many communication steps, so an infinite execution
cannot consist of either infinitely many communications or an infinite pure
reduction between them.

For uniqueness, pure confluence determines the owner's witness up to
conversion. Substituting convertible witnesses gives convertible continuation
types and terms. Apply the induction hypothesis to that continuation. At the
terminal, normalize the result record. The equations defining $\Run$ respect
pure conversion, so taking a pure reduction earlier or later gives the same
complete result.
\end{proof}

The execution lemma concerns one finite scripted connection. It supplies neither
termination of a recursive process calculus nor deadlock freedom for an
arbitrarily wired network of channels.

\subsubsection{Role symmetry}

Let $\pi:\Roles\cong\Roles$ be a permutation. It acts by
\begin{equation}
 \pi\End\langle R_i\rangle_i
   =\End\langle R_{\pi^{-1}(i)}\rangle_i,
 \qquad
 \pi(\Q_r(x:A).H)=\Q_{\pi(r)}(x:A).\pi H.
 \label{dep:eq:15}
\end{equation}
Data types and chosen values are unchanged. The corresponding endpoint
transport retains its data term and changes its role and global annotation.

\begin{lemma}[Endpoint transport]
\label{lem:role-transport}
$\Local{\pi(i)}{\pi G}\equivdef\Local{i}{G}$.
\end{lemma}
\begin{proof}
At a terminal, the inverse permutation in \eqref{dep:eq:15} selects $R_i$. At a choice,
$\pi(i)=\pi(r)$ iff $i=r$, so the same sum or product is selected in \eqref{dep:eq:4}.
Induct on the continuation.
\end{proof}

Define the transported profile by $(\pi\vecp)_j=p_{\pi^{-1}(j)}$, with its
annotation transported by \cref{lem:role-transport}. On execution records,
define $\pi_*(t,(z_i)_i)=(t,(z_{\pi^{-1}(j)})_j)$: trace values are retained
and the terminal record is permuted in the same direction.

\begin{lemma}[Equivariance]
\label{thm:equivariance}
For every well-typed profile,
\begin{equation}
 \Run_{\pi G}(\pi\vecp)=\pi_*\Run_G(\vecp).
 \label{dep:eq:16}
\end{equation}
Role transport also takes each communication step to a communication step.
\end{lemma}
\begin{proof}
The transported owner supplies the same pair and therefore the same witness.
The transported receivers apply the same functions to that witness. Induct on
the residual script. The terminal record is reordered exactly by $\pi_*$.
The same argument applied to \eqref{dep:eq:14} proves the reduction statement.
\end{proof}

Symmetry is a property of the family of roles and endpoints. An individual
endpoint judgment names the role it implements, just as an ordinary function
judgment names its arguments. There is no permanently privileged role in
\eqref{dep:eq:2}--\eqref{dep:eq:16}.
